
\begin{tabular}{p{\linewidth}}
  \footnotesize
  \justifying
    \cellcolor{cvbackgroundgray}
\noindent Mechanical Engineer and Machine Learning Specialist focused on physics-aware modeling and AI-based condition monitoring for industrial cyber-physical systems. Strong ability to interpret sensor behavior, analyze physical degradation mechanisms, and define physically consistent anomaly detection criteria before implementing scalable ML solutions. Experienced in high-frequency industrial time-series modeling, data pipeline automation, and end-to-end system deployment. Demonstrated industrial impact while pursuing mathematically grounded approaches to predictive maintenance and structured system modeling.
 My recent work increasingly emphasizes structured and topology-aware system representations that encode physical relationships and constraints, reflecting a growing research-oriented focus on interpretable and physically grounded AI models for complex engineering architectures.
\end{tabular}



%  data engineer 
% Machine Learning Engineer with a strong data engineering focus, experienced in designing scalable data architectures and automating high-frequency industrial time-series pipelines for AI-driven predictive maintenance systems. Skilled in building robust ETL workflows, integrating shop-floor IoT streams with cloud-based analytics platforms (Azure ML, Snowflake, Databricks, OPC UA, MQTT), and enabling reliable end-to-end ML deployment. Demonstrated measurable business impact through improvements in asset availability, production output, and operational efficiency. Mechanical engineering background supporting physics-aware system design.


 
% machine learning 
% Machine Learning Engineer and Mechanical Engineer specializing in AI-driven predictive maintenance and physics-aware condition monitoring for industrial systems. Experienced in developing, validating, and deploying ML/DL models for anomaly detection and RUL prediction, supported by scalable data engineering pipelines for high-frequency time-series data. Strong integration of IoT data streams across shop-floor and cloud analytics layers (Azure ML, Snowflake, Databricks, OPC UA, MQTT). Proven business impact through measurable improvements in asset availability and production performance.

% research
% Mechanical Engineer and Machine Learning Specialist focused on physics-aware modeling and AI-based condition monitoring for industrial cyber-physical systems. Strong ability to interpret sensor behavior, analyze physical degradation mechanisms, and define physically consistent anomaly detection criteria before implementing scalable ML solutions. Experienced in high-frequency industrial time-series modeling, data pipeline automation, and end-to-end system deployment. Demonstrated industrial impact while pursuing mathematically grounded approaches to predictive maintenance and structured system modeling.
% My recent work increasingly emphasizes structured and topology-aware system representations that encode physical relationships and constraints, reflecting a growing research-oriented focus on interpretable and physically grounded AI models for complex engineering architectures.

% \noindent Mechanical Engineer and Machine Learning researcher working at the intersection of artificial intelligence and engineering systems. My work focuses on physics-aware machine learning for complex cyber-physical systems, combining industrial domain knowledge with data-driven models to develop interpretable and reliable AI methods.
% In industrial settings, I have developed scalable machine learning pipelines for high-frequency sensor data, addressing anomaly detection, degradation modeling, and remaining useful life prediction for large-scale hydraulic systems. This experience has motivated my research interest in structured learning approaches that explicitly incorporate system topology, physical constraints, and engineering design knowledge into machine learning models.
% My long-term research goal is to develop machine learning methods that integrate physical structure, topology, and system architecture into learning algorithms, enabling robust AI systems capable of reasoning about complex engineered systems.